\section{Problem Formulation and LP Structure}
\label{sec:formulation}

\subsection{The Multi-Source CFLP}

Let $\J = \{1,\ldots,m\}$ be the set of potential facility locations and
$\I = \{1,\ldots,n\}$ the set of clients.
Each facility $j \in \J$ has a fixed opening cost $f_j > 0$ and a capacity
$q_j > 0$.
Each client $i \in \I$ has a demand $d_i > 0$.
The unit shipping cost from facility $j$ to client $i$ is $c_{ji} \geq 0$.

Let $y_j \in \{0,1\}$ indicate whether facility $j$ is open, and let
$x_{ji} \in [0, d_i]$ be the amount of client $i$'s demand served by
facility $j$.
The MSCFLP is:

\begin{align}
  \min \quad & \sum_{j \in \J} f_j y_j
               + \sum_{j \in \J} \sum_{i \in \I} c_{ji} x_{ji}
               \label{eq:obj} \\
  \text{s.t.} \quad
  & \sum_{j \in \J} x_{ji} = d_i
    && \forall i \in \I
    \label{eq:demand} \\
  & \sum_{i \in \I} x_{ji} \leq q_j y_j
    && \forall j \in \J
    \label{eq:capacity} \\
  & x_{ji} \leq d_i y_j
    && \forall j \in \J,\; i \in \I
    \label{eq:vub} \\
  & y_j \in \{0,1\}
    && \forall j \in \J
    \label{eq:binary} \\
  & x_{ji} \geq 0
    && \forall j \in \J,\; i \in \I.
    \label{eq:nonneg}
\end{align}

Constraints~\eqref{eq:demand} ensure all demand is satisfied.
Constraints~\eqref{eq:capacity} enforce facility capacities.
The variable-upper-bound (VUB) constraints~\eqref{eq:vub} prevent routing
demand to closed facilities.
Together, constraints~\eqref{eq:capacity} and~\eqref{eq:vub} ensure that
$y_j = 0$ implies $x_{ji} = 0$ for all $i$, tightening the LP relaxation
substantially compared to formulations without VUB constraints.

The LP relaxation is obtained by replacing~\eqref{eq:binary} with
$0 \leq y_j \leq 1$.
We denote its optimal solution by $(\xLP, \yLP)$ and its objective value
by $z^{\text{LP}}$.
Unless explicitly stated otherwise, gaps reported in this paper are measured
relative to $z^{\text{LP}}$:
\[
  \text{gap}(\%) = \frac{z^H - z^{\text{LP}}}{z^{\text{LP}}} \times 100,
\]
where $z^H$ is the heuristic objective.
This gives a common lower-bound-based measure across all algorithmic variants.
When comparing with published tables, we also report objective values because
published gaps may use different lower bounds or hardware.

\subsection{LP Solution Structure}
\label{sec:lp_structure}

The LP solution partitions facilities into three disjoint sets based on
$\yLP_j$:
\begin{itemize}
  \item $\Jone = \{j \in \J : \yLP_j = 1\}$: LP-open facilities,
  \item $\Jfrac = \{j \in \J : 0 < \yLP_j < 1\}$: LP-fractional
        facilities,
  \item $\Jzero = \{j \in \J : \yLP_j = 0\}$: LP-closed facilities.
\end{itemize}
The \emph{kernel} is $\K = \Jone \cup \Jfrac$, following~\citet{guastaroba2012kernel}.

\paragraph{Why the kernel predicts integer solutions.}
The VUB constraints~\eqref{eq:vub} force $x_{ji} \leq d_i y_j$: no demand
can be routed to a facility without opening it, even fractionally, in the LP.
A LP-closed facility ($\yLP_j = 0$) therefore carries no LP flow; the LP
found it suboptimal to open facility $j$ even fractionally.
This is a useful signal: a facility that is LP-closed must overcome both fixed
opening cost and capacity interactions before it becomes attractive in an
integer solution.
The computational evidence in \citet{guastaroba2012kernel} shows that the
LP-defined kernel captures most facilities used in high-quality solutions on the
large-scale benchmark families.
This motivates the Kernel Search strategy: rather than solving the full
$m$-facility MILP, start from $\K$ and test promising LP-closed facilities
through restricted subproblems.

\paragraph{Reduced costs.}
Let $\pi_i \geq 0$ be the dual variable for demand constraint~\eqref{eq:demand}
and $\alpha_j \leq 0$ the dual variable for capacity
constraint~\eqref{eq:capacity}.
For the LP relaxation, the \emph{reduced cost} of shipping variable $x_{ji}$
is:
\begin{equation}
  \rc{ji} = c_{ji} - \pi_i - \alpha_j.
  \label{eq:rc}
\end{equation}
At LP optimality, $\rc{ji} \geq 0$ for every variable at its lower bound
($x_{ji} = 0$) and $\rc{ji} = 0$ for every basic variable ($x_{ji} > 0$).
A variable with large positive reduced cost is far from being optimal to
include: the LP would need a much cheaper shipping cost, or a much better
dual price, before it would use that arc.
This signal guides both column generation (Section~\ref{sec:cg}) and
edge selection (Section~\ref{sec:edge_selection}).

\paragraph{Why column generation is natural.}
The LP solution of the strong MSCFLP formulation is often sparse relative to the
complete $|\J||\I|$ shipping-variable matrix.
Many shipping variables are zero at LP optimality and have nonnegative reduced
cost.
Column generation exploits this structure by leaving such variables outside the
restricted master problem until their reduced cost indicates that they may
improve the LP objective.
Study~\ref{sec:study_lp} reports the observed active-column counts and checks
LP agreement against the full formulation where the full LP can be solved.
